\chapter{Conclusões e Recomendações}
\label{cap:06}
\label{cap:conclusoes}

O objetivo geral deste trabalho foi desenvolver e avaliar o GUY, uma extensão para o VS Code com foco inicial na análise sintática de código Python e na geração do Grafo de Fluxo de Controle (GFC). A pesquisa foi concluída com a implementação e a validação técnica da ferramenta no escopo definido, contemplando a geração visual do GFC, o cálculo da Complexidade Ciclomática \(V(G)\), a navegação entre grafo e código-fonte e a execução local da análise. Assim, os experimentos verificaram as funcionalidades e os resultados métricos previstos para os exemplos avaliados; possíveis aplicações profissionais ou educacionais permanecem como hipóteses para estudos posteriores.

Na análise das ferramentas disponíveis, foram examinadas soluções voltadas à qualidade e à análise de código, como SonarQube, Codacy e ESLint. A comparação qualitativa considerou os recursos observados nessas ferramentas e na proposta do GUY, incluindo a forma de apresentação dos diagnósticos e a visualização do GFC. Esse levantamento ajudou a delimitar o foco da extensão em uma representação visual e interativa do fluxo de controle, sem estabelecer conclusões sobre superioridade ou adequação a um público específico.

O levantamento de requisitos definiu os requisitos funcionais e não funcionais da extensão, incluindo execução local, integração com o VS Code, geração do GFC, cálculo de \(V(G)\), indicação de caminhos independentes e interação entre grafo e código-fonte. Essa especificação, representada também pelos diagramas de casos de uso e de sequência, orientou a implementação e estabeleceu os critérios usados posteriormente na validação.

Durante a prototipagem e o design, foi concebido o fluxo de uso da ferramenta no editor, com abertura de uma Webview lateral e apresentação conjunta do GFC e da Complexidade Ciclomática. O protótipo organizou a navegação a partir de nós, arestas e caminhos independentes e traduziu os requisitos definidos em uma interface integrada ao VS Code. Sua usabilidade, porém, não foi avaliada com usuários.

Na implementação, a extensão foi desenvolvida em \textit{TypeScript}, com integração ao VS Code, análise sintática de Python por Tree-sitter em WebAssembly, construção do GFC, cálculo de métricas estruturais e renderização interativa em uma Webview construída com React e \texttt{@xyflow/react}. O protótipo funcional foi organizado em módulos responsáveis por análise, construção do grafo, métricas, navegação no editor e interface visual. Essa implementação permitiu testar as funcionalidades com exemplos concretos.

Na validação técnica, foram executados experimentos com três arquivos de exemplo em Python para verificar a coerência entre as estruturas de controle analisadas, os grafos gerados e as métricas calculadas. Foram registradas tabelas com nós, arestas, decisões, componentes, \(V(G)\) e caminhos independentes para cada experimento. Os resultados são válidos para os cenários avaliados e não constituem evidência de usabilidade, aprendizagem ou desempenho em projetos reais.

Os resultados quantitativos registrados no Capítulo \ref{cap:05} indicaram que, no Experimento I, o código com laço \texttt{while}, dois comandos \texttt{if}, \texttt{continue} e \texttt{break} gerou 13 nós, 15 arestas, 3 decisões, 1 componente, \(V(G)=4\) e 4 caminhos independentes. No Experimento II, o código condicional simples gerou 8 nós, 8 arestas, 1 decisão, 1 componente, \(V(G)=2\) e 2 caminhos independentes. No Experimento III, a função com laço \texttt{for}, condicional \texttt{if} e retorno gerou 9 nós, 10 arestas, 2 decisões, 1 componente, \(V(G)=3\) e 3 caminhos independentes. Esses valores confirmaram, nos exemplos analisados, a coerência entre a topologia do GFC e a fórmula da Complexidade Ciclomática.

Do ponto de vista qualitativo, a implementação incluiu integração ao VS Code, execução local da análise, visualização interativa do GFC, relacionamento entre elementos do grafo e trechos do código-fonte, modos de visualização simplificado e detalhado e análise de arquivo completo, seleção ou função atual. Esses recursos caracterizam possibilidades de uso analítico e, potencialmente, educacional, mas sua efetividade para esses fins não foi avaliada nesta pesquisa.

As limitações observadas também delimitaram o alcance da solução. A primeira limitação foi o foco da versão atual em Python e em estruturas de controle comuns, sem representação completa de exceções, chamadas entre funções, recursão, geradores, programação assíncrona e efeitos dinâmicos de execução. A segunda limitação decorreu da própria análise estática, que não substitui a execução de testes nem garante a identificação de todos os comportamentos possíveis em tempo de execução. A terceira limitação esteve relacionada à visualização de grafos muito grandes, pois programas com muitas estruturas aninhadas podem gerar representações visualmente densas, mesmo com avisos, modo simplificado e ocultação da listagem de caminhos independentes em grafos extensos.

\section{Trabalhos Futuros}

Em continuidade à primeira limitação, recomenda-se ampliar o analisador para cobrir mais recursos da linguagem Python, incluindo exceções, chamadas entre funções, recursão, geradores e programação assíncrona, preservando a coerência entre os elementos sintáticos reconhecidos e a construção do GFC.

Para tratar a segunda limitação, recomenda-se complementar a validação estática com testes automatizados baseados em exemplos de referência, de modo que os grafos e métricas gerados possam ser comparados continuamente com resultados esperados durante a evolução da ferramenta.

Quanto à terceira limitação, recomenda-se aprimorar os recursos de navegação e organização visual para grafos extensos, incluindo estratégias de agrupamento, filtragem ou exportação da visualização, a fim de reduzir a densidade visual em programas maiores. Também se recomenda realizar estudos com usuários e protocolos explícitos para avaliar usabilidade, aprendizagem e utilidade em contextos profissionais ou educacionais.
